% 这一章写综合使用 跨域一致性信号和竞争信号 来恢复部分 AES-128的密钥
% 我觉得要分为以下几个小节来写：
% 1.
% 2. 攻击的原理。为什么可以通过 AES T-Table 的访问模式来恢复密钥？这里需要分析 AES T-Table 实现的访存模式，以及我们测量到的信号和访存模式之间的关系。这一节还需要写为什么我们没办法在真实的机器中去完整恢复 AES 的全部密钥（需要结合已有的一些ch5下的实验，说明信号噪音高，空间精度不够，探测窗口不够，没办法通过这两类信号清楚定位到某一缓存行是否被访问）。导致只能使用T-table的不对齐性（可以参考Cohere+Reload论文，但是不要提及）来恢复部分密钥，这部分的原理也要讲，可以举例当前机器上AES结构（aes_base=0x409000 te0_va=0x409f60 te0_file_off=0xf60 te0_gpa=0x101c74f60 te0_page_gpa=0x101c74000，可画图来展现这个对齐结构）。
% 3. 攻击系统的总体设计。攻击系统分为两个阶段：第一阶段是确定T-table所在的页和页内偏移（脚本report/run_53_seed_range_v2.py，实验结果report/ch5/exp5_3_seed_success等），第二阶段是基于第一阶段的结果来构建攻击模板并进行密钥恢复。第一阶段能获取te0_gpa和te0_page_gpa，然后根据这些信息来定位目标页和页内偏移。第二阶段则根据te0_gpa和te0_page_gpa，得到T-table的不对齐的具体情况，然后构建攻击模板来针对特定的密钥字节进行恢复。（这一节可以先大体介绍总体的流程，后续再在后续章介绍具体的流程）
% 3.1 定位T-table所在的页和页内偏移（分析脚本来写具体的攻击流程）。注意这个攻击并不能精准一次获得目标页。而是得到大概的一个目标页的簇。注意这里为了让真实目标簇最后的分数高，除了Coheren信号、NPT访问频率外，还根据实际情况引入了一些其他特定的特征（分析代码来帮我写出来）

% 4. 真实实验
% 4.1 实验平台与环境：介绍硬件平台、软件环境（包括每次测量 host可以发请求控制 guest执行加密行为）、AES 实现细节（这一段需要可以结合openssl源码来详细分析）等
% 4.2 AES T-Table 物理页定位：结合实际实验过程介绍第一阶段的攻击流程，分析实验结果，说明最终是如何定位到目标页的。注意要写最终定位可以将目标页从候选页簇中区分出来（前三的），说明引入的特征是有效的。需要介绍具体我们筛选的特征。这一节需要结合具体的一些实验数据来说明（不要与前面系统设计重复了）。可以说也不是百分百都能目标页到rank 1,但是能非常有效地将目标页从候选页簇中区分出来（前rank3）。
% 4.3 AES T-Table 部分密钥恢复：结合实际实验过程介绍第二阶段的攻击流程，分析实验结果，说明最终是如何通过构建的模板来恢复部分密钥的。需要介绍具体我们构建的特征，以及为什么这些特征能够区分正确的密钥候选值和错误的密钥候选值。需要结合具体的一些实验数据来说明（不要与前面系统设计重复了）。可以说也不是百分百都能将正确密钥候选值排在第一位，但是能非常有效地将正确密钥候选值从错误密钥候选值中区分出来（前rank10）。

\chapter{AES-128 的部分密钥恢复攻击}

第~\ref{chap:channel}~章对 SEV-SNP 环境下的两类侧信道信号进行了量化表征。跨域一致性逐出信号的 ROC-AUC 为 0.9351，最优阈值处 $P_D = 1.0000$、$P_{FA} = 0.0001$，信道容量为 0.9993 比特/次探测；物理内存竞争信号的 ROC-AUC 为 0.9007，最优阈值处 $P_D = 0.8709$、$P_{FA} = 0.0552$，信道容量为 0.5652 比特/次探测。将上述信号应用于真实 AES 加密时，存在两个制约因素：其一，一致性信号的有效感知范围以 2KB 为粒度，无法区分同一块内不同缓存行是否被访问；其二，竞争信号在真实 AES 多轮访问背景下信噪比下降，难以对单条缓存行做出可靠判决。在当前实验平台的信号条件下，逐缓存行的完整密钥恢复不可行，需要利用 T-Table 的内存不对齐特性，将攻击目标转化为可通过页面级观测区分的问题。本章围绕这一思路构建两阶段攻击流程：第一阶段利用 NPT 访问位差分扫描定位 \texttt{Te0} 所在的物理页，第二阶段在已定位的页上采集内存竞争信号序列，以模板匹配方法恢复部分密钥字节。

\section{真实实验}
\label{sec:real_exp}

\subsection{实验平台与环境}
\label{subsec:exp_platform}

\subsubsection{硬件平台与测量配置}

本章实验的硬件平台与第~\ref{chap:channel}~章相同，系统配置见表~\ref{tab:exp42_setup}所示。宿主机探测线程通过 \texttt{sched\_setaffinity} 与 QEMU vCPU 线程绑定在同一物理核的两个超线程上，使两者共享 L1 和 L2 缓存，以保证一致性信号的观测强度。宿主机侧利用 KVM 内核扩展接口与 \texttt{Te0} 目标页交互，该接口提供两类操作，1. 将来宾物理地址（GPA）解析为宿主物理地址（HPA），用于后续的缓存行延迟测量；2. 批量读取并清除 NPT 页表项中的 Accessed 位，用于第一阶段的页面定位。宿主机通过 Unix 套接字与内核扩展通信，每条指令的往返时延约为数十微秒，不影响毫秒级的信号采集流程。

\subsubsection{来宾虚拟机配置}

来宾虚拟机为单 vCPU 的 SEV-SNP 加密模式，由 AMDSEV 上游分支维护的 QEMU/KVM 启动。来宾内核以 \texttt{nokaslr norandmaps} 选项启动，禁用内核和用户态地址随机化，使 \texttt{Te0} 在二进制文件中的页内偏移在每次启动后保持一致。用户态运行 BusyBox 环境，主体为静态链接的目标程序，内嵌 OpenSSL 3.4.0 的 \texttt{no-asm} 编译版本。

来宾程序在启动后监听 TCP 9000 端口，对收到的每条 16 字节明文执行一次 AES-128 加密并返回密文。宿主机通过向该端口发送随机明文来触发加密操作，在 NPT 扫描阶段每轮发送若干请求，在信号采集阶段按受控时序逐一触发。这种网络服务接口是攻击者唯一需要的外部控制手段，不依赖对来宾内存或执行流的任何直接干预。

\subsubsection{AES 实现路径验证}

OpenSSL 3.4.0 默认会在 x86-64 平台上使用汇编优化版本（\texttt{aesni\_encrypt}），该路径通过 \texttt{aesenc} 硬件指令完成加密，不产生可观测的内存查表访问。为使加密退回到软件 T-Table 实现，目标程序在编译时使用 \texttt{no-asm} 标志，同时在 QEMU CPU 模型中移除 \texttt{aes} 特性位，使来宾 CPUID 不报告 AES-NI 支持。两项配置共同作用，确保运行时进入 \texttt{AES\_encrypt} 函数，依赖 \texttt{Te0}--\texttt{Te3} 四张软件查表完成每轮运算。

路径验证结果如表~\ref{tab:aes_path_annotate} 所示。在 \texttt{no-asm} 配置下，\texttt{perf annotate} 显示 \texttt{Te0}-\texttt{Te3} 相关变址内存读取占热点指令的 7.48\%；默认 \texttt{asm} 配置下该比例为 0.00\%，热点集中于 \texttt{\_x86\_64\_AES\_encrypt\_compact} 函数中的 \texttt{vaesenc} 系列指令。两种配置的访存特征差异明显，确认 \texttt{no-asm} 版本产生的 T-Table 查表行为可供侧信道观测。

\begin{table}[htbp]
  \centering
  \caption{两种 AES 编译配置的 \texttt{perf annotate} 指令热点对比}
  \label{tab:aes_path_annotate}
  \begin{tabular}{lcc}
    \toprule
    指标 & \texttt{no-asm} T-Table & 默认 \texttt{asm} 紧凑路径 \\
    \midrule
    热点函数 & \texttt{AES\_encrypt} & \texttt{\_x86\_64\_AES\_encrypt\_compact} \\
    Te0-Te3 变址内存读取占比 & 7.48\% & 0.00\% \\
    \texttt{vaesenc} 占比 & 0.00\% & 62.14\% \\
    \bottomrule
  \end{tabular}
\end{table}

\section{AES T-Table 侧信道攻击原理}
\label{sec:aes_attack_principle}

\subsection{T-Table 实现的访存模式}

OpenSSL 在 x86-64 平台默认使用 AES-NI 硬件指令完成加密，该路径不产生可观测的内存查表行为。当编译时指定 \texttt{no-asm} 标志且 CPU 不报告 AES-NI 特性时，OpenSSL 退回到软件实现路径，调用函数 \texttt{AES\_encrypt}。该函数使用四张预计算查找表（Te0--Te3），每张表由 256 个 32 位整数组成，占连续 1024 字节内存空间。T-Table 将 SubBytes、ShiftRows 和 MixColumns 三步运算合并为单次内存读取，以提高执行效率，但也因此引入了与密钥相关的内存访问模式。

在明文与初始轮密钥异或后，第一轮对四个状态字的计算展开为

\begin{equation}
  \begin{aligned}
    t_0 &= \mathrm{Te0}[s_0 \gg 24] \\
    &\oplus \mathrm{Te1}[(s_1 \gg 16)\,\&\,\mathtt{0xff}] \\
    &\oplus \mathrm{Te2}[(s_2 \gg 8)\,\&\,\mathtt{0xff}] \\
    &\oplus \mathrm{Te3}[s_3\,\&\,\mathtt{0xff}] \\
    &\oplus \mathit{rk}[4]
  \end{aligned}
  \label{eq:aes_round1}
\end{equation}

其余三个输出字 $t_1, t_2, t_3$ 具有对称结构。状态字 $s_0$ 由明文字节与初始轮密钥逐字节异或构成

\begin{equation}
  s_0 = \bigl(p_0 \ \| \ p_1 \ \| \ p_2 \ \| \ p_3\bigr) \oplus \mathit{rk}[0]
  \label{eq:state_word}
\end{equation}

因此，\texttt{Te0} 的查表索引化简为 $e_0 = s_0 \gg 24 = p_0 \oplus k_0$，其中 $p_0$ 为已知明文字节，$k_0$ 为待求密钥字节。查表发生时，CPU 将包含 \texttt{Te0[$e_0$]} 的缓存行（64 字节）从 DRAM 载入缓存。由于每条缓存行恰好容纳 16 个连续的 32 位表项，整张 Te0 共分布在 16 条缓存行上，第 $j$ 条缓存行（$j = 0, \ldots, 15$）覆盖索引区间 $[16j,\ 16j+15]$。因此哪一条缓存行被载入取决于 $e_0 = p_0 \oplus k_0$ 的值，即 $e_0$ 唯一确定了本次加密触发的缓存行编号。若攻击者能观察到这次缓存行载入事件，便可推断 $e_0$ 所在的 16-索引区间，进而结合已知的 $p_0$ 推导出密钥字节 $k_0$。

AES-128 第一轮共产生 16 次 T-Table 查表访存，字节位置 0、4、8、12 访问 \texttt{Te0}，1、5、9、13 访问 \texttt{Te1}，2、6、10、14 访问 \texttt{Te2}，3、7、11、15 访问 \texttt{Te3}。后续轮次经 ShiftRows 和 MixColumns 扩散后，访存地址与密钥的直接对应关系消失，实际攻击以第一轮查表为目标。

\subsection{信号与访存模式的关系}

跨域一致性逐出信号的产生与来宾的缓存访问直接相关，当来宾执行 \texttt{Te0[$e_0$]} 查表时，CPU 将目标缓存行以来宾私有密钥（C-bit 置 1）从 DRAM 解密后载入缓存，此时缓存中存放的是来宾侧的明文内容。SEV-SNP 的双视图内存模型使宿主机通过同一物理地址的宿主侧映射（C-bit 置 0）看到的是不同内容，两者在缓存内的标签不一致，触发跨域一致性协议，宿主机持有的缓存副本随之被逐出失效。宿主机在来宾查表后再次探测同一物理地址时，须从 DRAM 重新读取，访问延迟明显高于缓存命中情况，形成 $H_1$ 状态。若来宾在测量窗口内未访问该缓存行，宿主机的缓存副本保持有效，探测延迟维持基线 $H_0$ 水平。

具体到 \texttt{Te0} 的访存模式，每条 64 字节缓存行覆盖 16 个连续表项。当查表索引 $e_0 = p_0 \oplus k_0$ 落在某条缓存行所覆盖的区间内，宿主机在该缓存行上的探测产生 $H_1$ 事件，否则维持 $H_0$。攻击者通过控制明文字节 $p_0$ 改变 $e_0$ 的取值，可以选择性地使 $e_0$ 落在或偏离指定缓存行所覆盖的 16 个索引范围，从而在该缓存行上构造出与密钥字节 $k_0$ 相关的 $H_1/H_0$ 信号序列。

物理内存竞争信号的形成机制与一致性逐出不同，但同样与 Te0 的查表访存直接关联。来宾执行 Te0[$e_0$] 查表时，目标缓存行若不在缓存中，CPU 会向 DRAM 发起一次读请求，将该 64 字节缓存行从内存控制器载入。若此时宿主机以不可缓存（NOCACHE）映射方式访问同一物理地址，两次内存访问会在 DRAM 总线层面产生争用，宿主机的读延迟因此上升，形成可观测的竞争事件。NOCACHE 路径不经过缓存，每次探测均直接到达内存控制器，因此不受来宾 C-bit 加密标签的影响，宿主机无需解密即可触发总线争用并测量延迟。

竞争信号的空间粒度由 DRAM 行粒度决定，理论上可达到 64 字节缓存行的精度，因为只有在来宾实际访问目标缓存行时，对该地址的 NOCACHE 探测才会产生延迟升高。一致性逐出信号虽然能够稳定产生，但空间分辨率较低；竞争信号在页面内定位精度较高，但其产生需要宿主机的探测与来宾的 DRAM 读操作在同一短暂时间窗口内重叠。对 Te0 而言，每次加密仅在查表发生的数纳秒内触发 DRAM 读事件，宿主机 NOCACHE 探测若未在该窗口内到达内存控制器，则不会观测到竞争延迟，信号的产生概率低于一致性逐出信号。

\subsection{完整密钥恢复的限制}
\label{subsec:attack_limits}

上述理论分析基于理想假设，即每次加密仅访问一条缓存行且宿主机可将观测精确定位到该缓存行。在实际 SEV-SNP 环境中存在下面三种制约原因，使逐缓存行的完整密钥恢复难以直接实现。

\begin{enumerate}
  \item 一致性信号的空间粒度为 2KB。第~\ref{chap:channel}~章的实验表明，当来宾访问某一 2KB 对齐范围内的任意缓存行时，宿主侧在该范围内任意位置的探测均可产生显著的 $H_1$ 延迟；受 DRAM Interleaving 机制影响，该范围内不同缓存行的信号强弱存在差异。对 \texttt{Te0} 而言，1024 字节共 16 条缓存行，完整表横跨约 1.5 个 2KB 窗口。宿主侧在某条缓存行上观测到的 $H_1$ 事件只能表明对应 2KB 范围内存在访存活动，无法进一步区分具体是哪条 64B 缓存行被访问，也就无法确定 $e_0$ 所落在的 16-索引区间。

  \item 竞争信号的时序要求较严。竞争信号的空间精度为 64B，原则上可区分单条缓存行，但其产生需要宿主机的 NOCACHE 探测与来宾的 DRAM 访问在时间上重叠。AES 第一轮在约数 10 ns内完成，四张表上的 16 次查表访存相互交织，宿主侧若要将探测窗口对准某一次特定表项的访问，需要较高的时序精度。第~\ref{chap:channel}~章的测量结果显示竞争信号的 $P_D = 0.8709$、$P_{FA} = 0.0552$，在实际 AES 执行的多次访问背景下，漏检率和误报率相比孤立探测时均有所上升，单独依靠竞争信号进行密钥恢复所需的观测轮数较多。

  \item  AES 加密的有效观测窗口较窄。攻击者只能通过 TCP 请求触发加密操作，但 AES 第一轮本身仅持续约数 10 ns，之后 AES操作的 ShiftRows 和 MixColumns 的扩散使后续轮次的访存与密钥失去直接对应关系，宿主机能够加以利用的时间窗口有限。
\end{enumerate}

上述三类因素共同表明，在当前实验平台的信号条件下，逐缓存行追踪 $e_0$ 取值在实践中面临较大困难。

\subsection{T-Table 不对齐结构与部分密钥恢复}
\label{subsec:te0_misalign}

上述困难的根源在于真实信号的精度不匹配，无论是一致性信号还是竞争信号，其在真实机器上的信号粒度不足以区分单次查表究竟命中了哪条缓存行。若 \texttt{Te0} 的全部 256 个表项均位于同一物理页，则无论密钥字节取何值，每次加密都会访问该页面，页面级观测无法提供任何推断依据。然而，\texttt{Te0} 在二进制文件中的起始偏移通常不在 4KB 页边界上，这使得 \texttt{Te0} 天然跨越两个连续物理页，不同索引区间的查表访存分布在两个不同的物理页上，将原来的缓存行观察问题转换成页面级观察问题。

以本实验使用的 OpenSSL 3.4.0 二进制为例，\texttt{Te0} 的页内起始偏移为 \texttt{0xF60}，对应来宾物理地址 \texttt{0x101C74F60}，其所在页（页 0）在 \texttt{0x101C74FFF} 处结束，剩余部分位于下一物理页（页 1）。页 0 末尾存放索引 0 至 39 共 160 字节的表项，页 1 开头存放索引 40 至 255 共 864 字节的表项。该布局如图~\ref{fig:te0_layout} 所示。

\begin{figure}[htbp]
  \centering
  \includegraphics[width=\textwidth]{pic/te0_layout.pdf}
  \caption{\texttt{Te0} 跨页内存布局示意（te0\_inpage\_off = 0xF60）}
  \label{fig:te0_layout}
\end{figure}

该布局使一次加密的查表行为产生可区分的页面访问模式。当 $e_0 \in [0, 39]$ 时，查表操作访问页 0；当 $e_0 \in [40, 255]$ 时，访问页 1。宿主机对两个物理页分别进行一致性信号探测，即可判断本次加密的查表索引属于哪个区间。对已知明文字节 $p_0$，由 $k_0 = p_0 \oplus e_0$ 可得密钥字节的候选集：若判决为页 0，则候选集为 $\{p_0 \oplus e \mid e \in [0,39]\}$，共 40 个候选值；若判决为页 1，则共 216 个候选值。两种情况分别将候选集从 256 缩减至约 15.6\% 或 84.4\%。

一致性信号的页面级判决将密钥字节候选集从 256 缩减到两个区间，但单次的粗粒度二值判决对密钥字节的直接约束有限，且一致性信号在高频重复探测中噪声较大，独立作为模板特征时区分能力不足。实验中对比两类信号后发现，内存竞争信号对正确密钥候选值的分类贡献更为显著：宿主机在每次加密结束后立即以不可缓存映射对 Te0 所在物理页的目标缓存行发起一组突发探测，记录各缓存行的延迟读数序列；当索引 $e_0$ 命中特定缓存行时，该缓存行在探测期间触发 DRAM 访问，延迟偏高；未命中时延迟偏低。不同索引值会使 $e_0$ 落在 Te0 内不同位置的缓存行，从而在各探测缓存行上产生可区分的计数分布。

积累多轮观测后，可以对每个候选密钥值 $k$ 根据竞争信号特征统计其吻合程度：若候选值 $k$ 等于真实密钥字节，则它所隐含的索引映射 $e_0 = p_0 \oplus k$ 与每轮实际的查表行为一致，该索引对应的内存竞争信号模型与观测特征匹配，累积得分偏高；若候选值 $k$ 错误，索引映射偏离真实访问位置，特征分布不匹配，得分偏低。基于此，密钥恢复采用高斯混合模型模板方法对 256 个候选密钥值逐一评分：离线阶段以已知密钥采集各索引值对应的内存竞争信号特征向量，为每个索引 $e$ 拟合条件概率模型 $\mathrm{GMM}[e]$；在线阶段对未知密钥下采集的特征向量序列 $\{\mathbf{x}_i\}$ 计算各候选值 $k$ 的累积对数似然

\begin{equation*}
  S(k) = \sum_{i} \log p\!\left(\mathbf{x}_i \mid \mathrm{GMM}[p_{0,i} \oplus k]\right),
\end{equation*}

取 $S(k)$ 最高的候选值作为密钥字节估计 $\hat{k}_0$。模板方法的具体实现在后续章节详述。

\section{攻击系统总体设计}
\label{sec:attack_overview}

基于上述分析，本文设计了一个两阶段的 AES-128 密钥恢复攻击系统。整个系统运行在宿主机侧，不修改来宾代码，不破解 SEV-SNP 内存加密，仅依赖宿主机对 NPT 的管理权限、跨域一致性逐出信号（用于 Te0 物理页定位与验证）和物理内存竞争信号（用于模板特征提取与密钥恢复）。图~\ref{fig:attack_pipeline} 展示了攻击系统的总体架构。系统由三个层次组成：位于上方的两个阶段框分别负责 Te0 物理页定位和密钥恢复，中间的 KVM 内核扩展节点提供硬件交互接口，底部的机密虚拟机（SEV-SNP）为攻击目标。

\begin{figure}[htbp]
  \centering
  \includegraphics[width=\textwidth]{pic/attack_pipeline.pdf}
  \caption{AES-128 密钥恢复攻击系统两阶段架构}
  \label{fig:attack_pipeline}
\end{figure}

攻击两阶段在功能上相互解耦，通过中间结果 \texttt{te0\_page\_gpa} 和 \texttt{te0\_inpage\_off} 串联。两个阶段均经由 KVM 内核扩展（通过 \texttt{nptctl.sock} 套接字）与硬件交互，宿主机用户态进程可在不修改来宾内存映射的前提下完成完整攻击流程。

\textbf{Te0 物理页定位：} 攻击者事先可从本地同版本二进制的符号表中获取 \texttt{Te0} 的页内偏移，但不知道其运行时 GPA。第一阶段通过 NPT Accessed 位的差分扫描来定位目标页面：在基线状态下扫描来宾 GPA 空间中被访问的页面，再在触发若干次 AES 加密后重新扫描，按 $\mathrm{score} = \mathrm{trigger\_hits} - \mathrm{baseline\_hits}$ 计算每个候选页面的差分得分，得分较高的页面簇为与加密操作相关的候选集。由于 AES 执行期间还会访问代码页、栈页等无关页面，差分评分存在多个高分候选，因此第一阶段引入跨域一致性逐出信号对候选页进行验证，对每个候选页计算触发加密前后的缓存行延迟差 $\Delta$ 及信号信噪比（SNR），结合差分得分构造综合评分。在实验中，真实 \texttt{Te0} 所在页的综合评分（804.2）高于其余候选页（次高约 71.5），将目标页从候选集中区分出来。第一阶段输出的 te0\_page\_gpa 为 Te0 所在物理页的基地址，te0\_inpage\_off 为 Te0 在该页内的字节偏移，后者决定了 Te0[0..39] 与 Te0[40..255] 的分界位置。

\textbf{模板构建与密钥恢复：} 获得 \texttt{te0\_page\_gpa} 和 \texttt{te0\_inpage\_off} 后，攻击者可计算其 T-Table 不对齐边界，即 T-Table 跨页的边界。第二阶段首先进行边界验证，通过控制明文使索引在页面边界附近变化，观察两个页面的一致性信号是否在边界处发生翻转，以此确认 $e^\ast$ 的取值。随后进入离线模板构建步骤：攻击者以已知密钥对字节位置 0 反复触发加密，使 $e_0 = p_0 \oplus k_0$ 遍历 0 至 255，对每个 $e_0$ 采集多维内存竞争信号特征向量，并对该索引值的样本拟合含 1 至 2 个分量的对角协方差高斯混合模型，建立包含 256 个条目的离线模板库。在线密钥恢复步骤中，攻击者以未知密钥触发若干次加密并采集特征向量序列，对 256 个候选密钥字节 $k$ 依式~\ref{eq:gmm_score} 分别计算累积对数似然，取得分最高者为密钥字节估计 $\hat{k}_0$。上述流程对字节位置 0、4、8、12 分别执行，覆盖 AES 第一轮对 \texttt{Te0} 的全部四次查表。

\subsection{T-Table 物理页定位}
\label{sec:phase1}

该阶段需要确定两个量，Te0 所在物理页的基地址 $\mathit{te0\_page\_gpa}$ 和 Te0 在该页内的字节偏移 $\mathit{te0\_inpage\_off}$。两者共同定义了 Te0 跨页的分界位置，供第二阶段计算不对齐边界 $e^\ast$ 使用。其中 $\mathit{te0\_inpage\_off}$ 可在攻击发起前离线提取：从本地同版本二进制文件读取 Te0 的符号地址，结合其所在 ELF LOAD 段计算文件偏移，再对页大小取模，得
\begin{equation}
  \mathit{off} = \mathit{FileOffset}(\mathrm{Te0}) \bmod 4096,
  \label{eq:inpage_off}
\end{equation}
由此推导 $e^\ast = \lfloor (4096 - \mathit{off}) / 4 \rfloor - 1$。只要宿主机与来宾运行同一版本的目标程序，$\mathit{off}$ 在每次来宾启动后保持不变。运行时需要确定的是 $\mathit{te0\_page\_gpa}$，即 Te0 在来宾 GPA 空间中的物理页基地址。

确定 $\mathit{te0\_page\_gpa}$ 的基本手段是 NPT Accessed 位差分扫描。AMD NPT 为每个页表项维护一个硬件 Accessed 位，来宾 vCPU 访问对应物理页时由 MMU 自动置位，宿主机通过 KVM 内核扩展可批量清除这些位并在一段时间窗口后重新扫描，从而获知哪些页面在该窗口内被来宾访问过。每轮扫描分基线与触发两个阶段进行：基线阶段清除 Accessed 位后等待 $T$ 毫秒，使内核心跳、栈操作等背景活动积累到各页的 $\mathrm{base}$ 计数中；触发阶段再次清除后立即向来宾 AES 服务发送 $N$ 次随机明文请求，将窗口聚焦于加密执行期间的实际访存，计入各页的 $\mathrm{trig}$ 计数。扫描范围覆盖来宾两段 GPA 区域：低地址段等于虚拟机配置的物理内存大小，高地址段为 SEV-SNP 私有页面在宿主 GPA 视图中 4 GB 以上的别名区，两段均需覆盖以防目标页落在任一区域。经 $R$ 轮重复，差分得分定义为
\begin{equation}
  \mathrm{score}(p) = \mathrm{trig}(p) - \mathrm{base}(p).
  \label{eq:npt_score}
\end{equation}
算法~\ref{alg:npt_diff_scan} 给出上述过程的完整描述。

\begin{algorithm}[htbp]
  \SetAlgoLined
  \caption{NPT Accessed 位差分扫描}
  \label{alg:npt_diff_scan}
  \KwIn{迭代轮数 $R$；每轮 AES 请求数 $N$；基线等待时间 $T$ ms}
  \KwOut{候选页列表 $\mathcal{L}$，按差分得分降序排列}
  对所有页面 $p$，初始化 $\mathrm{base}[p] \leftarrow 0$，$\mathrm{trig}[p] \leftarrow 0$\;
  \For{$r \leftarrow 1$ \KwTo $R$}{
    \tcp{基线阶段：采集来宾背景访问}
    \textsc{NptClear}()\;
    等待 $T$ 毫秒\;
    $\mathcal{B} \leftarrow \textsc{NptScan}()$\;
    对所有 $p \in \mathcal{B}$，令 $\mathrm{base}[p] \mathrel{+}= 1$\;
    \BlankLine
    \tcp{触发阶段：采集 AES 加密期间的访问}
    \textsc{NptClear}()\;
    向来宾 AES 服务发送 $N$ 次随机明文请求\;
    $\mathcal{T} \leftarrow \textsc{NptScan}()$\;
    对所有 $p \in \mathcal{T}$，令 $\mathrm{trig}[p] \mathrel{+}= 1$\;
  }
  \BlankLine
  对所有页面 $p$，令 $\mathrm{score}(p) \leftarrow \mathrm{trig}[p] - \mathrm{base}[p]$\;
  \KwRet{$\mathcal{L}$，按 $\bigl(\mathrm{score}(p),\; \mathrm{trig}[p],\; {-\mathrm{base}[p]}\bigr)$ 三元组降序排列}
\end{algorithm}

在差分评分得到候选列表后，将依次经过评分筛选和页簇聚类两步处理，将范围收窄到 Te0 所在区域附近。评分筛选保留需要同时满足 $\mathrm{score}(p) \geq \theta$ 与 $\mathrm{trig}(p) > \mathrm{base}(p)$ 的页面。前一条件剔除与 AES 执行弱相关的页面；而后一条件过滤背景态本就高频访问的常驻页，此类页面两次计数相近，$\mathrm{trig}$ 不高于 $\mathrm{base}$，即便 $\mathrm{score}$ 偶尔达到阈值也会被排除。Te0--Te3 四张表在每次 AES 加密中均会被访问，相应物理页的 $\mathrm{trig}$ 接近轮数上限 $R$、$\mathrm{base}$ 接近零，两个条件均能满足。随后将物理地址相邻且间距不超过 $\delta$ 页的候选归为同一簇，按下式三元组字典序排名，取最优簇 $\mathcal{C}^*$：
\begin{equation}
  f(\mathcal{C}) = \Bigl(\,\textstyle\sum_{p \in \mathcal{C}} \mathrm{score}(p),\quad -|\mathcal{C}|,\quad {-\mathrm{span}(\mathcal{C})}\,\Bigr).
  \label{eq:cluster_score}
\end{equation}
三个分量分别对应簇的总差分得分、页面数和地址跨度，以得分更高、页数更少、跨度更小为排序依据。Te0--Te3 在全局数据段中连续排列，所在物理页总分高、跨度小，在排名中靠前。

最优簇 $\mathcal{C}^*$ 确定后，结合离线得到的 $\mathit{off}$，在簇内筛选出满足页内偏移条件的候选页，通常剩余一到两个。此时以跨域一致性逐出信号进行最终确认：控制明文令 Te0 的查表索引落在候选页所覆盖的表项范围内，触发加密，在宿主侧测量各候选物理页的缓存行平均延迟。存放 Te0 的真实物理页在来宾查表后产生跨域逐出，宿主侧延迟升高幅度约为数百时钟周期；误入候选的无关页面不触发相应查表行为，逐出幅度小，延迟差值 $\Delta$ 偏低。取 $\Delta$ 较大者确认为 $\mathit{te0\_page\_gpa}$，连同 $\mathit{te0\_inpage\_off}$ 传递至第二阶段。

\subsection{模板构建与密钥恢复}
\label{sec:phase2}

第二阶段以第一阶段确定的 $\mathit{te0\_page\_gpa}$ 和 $\mathit{te0\_inpage\_off}$ 为输入，分为离线模板构建和在线密钥恢复两部分。离线阶段在已知密钥的受控环境中为 256 个查表索引值分别建立侧信道特征的概率模型；在线阶段对未知密钥下采集的观测序列逐一打分，取累积得分最高的候选值作为密钥估计。

模板构建的目标是为每个查表索引值 $e \in [0, 255]$ 建立一个条件概率模型 $\mathrm{GMM}[e]$，描述当来宾以索引 $e$ 访问 Te0 时，宿主机所能观测到的侧信道特征向量分布。在已知密钥 $k_b^{\mathrm{known}}$ 的条件下，攻击者遍历 $e = 0, 1, \ldots, 255$，将明文第 $b$ 字节设为 $p_b = e \oplus k_b^{\mathrm{known}}$，迫使来宾每次加密以指定索引 $e$ 查表。每次加密结束后，宿主机立即对特征缓存行集合 $\mathcal{R}$ 中的每条缓存行发起一组内存竞争突发探测，记录各行的延迟读数序列作为一条原始观测，对每个 $e$ 重复采集 $S$ 条。

特征缓存行集合 $\mathcal{R}$ 的选取考虑两类来源。第一类是跨页边界附近的缓存行：这些缓存行覆盖的索引区间横跨 $e^*$，当索引从边界一侧切换到另一侧时，对应物理页的访存状态发生翻转，信号变化幅度较大，区分能力较强。第二类是对 Te0 所在物理页执行缓存行延迟扫描后，按信号强度排名靠前的若干缓存行，用于覆盖边界附近缓存行无法反映的同表热点区域。

每条原始观测包含 $D = |\mathcal{R}|$ 条缓存行的 Contention 信号延迟序列。为降低单次读数受调度抖动干扰的影响，将每条缓存行 $j$ 的序列压缩为一个整数计数值 $x_j$，即延迟读数落在预设区间 $[d_j^{\mathrm{lo}},\, d_j^{\mathrm{hi}})$ 内的次数：

\begin{equation}
  x_j = \#\bigl\{\, t \in \text{burst}_j \;\big|\; d_j^{\mathrm{lo}} \leq t < d_j^{\mathrm{hi}} \,\bigr\}, \quad j = 1, \ldots, D
  \label{eq:feature_component}
\end{equation}

突发探测内多次读数的计数相当于对随机噪声进行平均，使同一索引在不同采集轮次下的特征值更为集中。每条缓存行的区间端点 $[d_j^{\mathrm{lo}},\, d_j^{\mathrm{hi}})$ 通过对训练数据的区间扫描独立确定：遍历候选下界，对每个候选区间计算该缓存行在不同索引样本之间的 Fisher 准则分数，即组间方差与组内方差之比，取分数最高的区间作为最终的延迟区间。该过程在每条缓存行上独立执行，使每个特征维度的区分能力达到最大。

对同一索引 $e$ 采集的 $S$ 条特征向量，其分布通常呈现双态特征：当目标数据在探测时仍驻留于缓存时，各维度计数偏低；当数据已被逐出、探测触发 DRAM 访问时，各维度计数普遍偏高。由于同步精度的限制，同一索引下两种状态均可能出现，实测样本中两种状态的计数值各自接近正态分布，整体呈现双峰形态。高斯混合模型能够以较少的参数描述这种混合分布，其均值、方差和权重等参数均有明确的物理对应，因此本文对每个索引 $e$ 拟合一个包含两个分量的对角高斯混合模型：

\begin{equation}
  p(\mathbf{x} \mid e) = \sum_{c=1}^{2} w_c^{(e)}\,
  \mathcal{N}\!\left(\mathbf{x};\; \boldsymbol{\mu}_c^{(e)},\; \mathrm{diag}\!\left(\boldsymbol{\sigma}_c^{(e)\,2}\right)\right)
  \label{eq:gmm_model}
\end{equation}

其中 $w_c^{(e)}$ 为混合权重，$\boldsymbol{\mu}_c^{(e)}$ 和 $\boldsymbol{\sigma}_c^{(e)\,2}$ 分别为第 $c$ 个分量的均值向量和方差向量。协方差取对角形式，假设各维度在给定分量下条件独立，在减少待估参数数量的同时降低有限样本下的过拟合风险。

模型参数通过期望最大化算法从采集到的样本中估计。该算法交替执行 E 步和 M 步：E 步根据当前模型参数计算每个样本属于各分量的后验概率，M 步利用这一软分配结果对各分量的均值、方差和权重进行加权更新。两步交替迭代至参数收敛，对数似然在迭代过程中单调不减。由于该算法对初始值较为敏感，本文采用基于物理先验的初始化策略，即对每个训练样本的各维度计数求和得到标量总计数，以总计数最小的样本均值初始化缓存命中分量，以总计数最大的样本均值初始化 DRAM 访问分量，使两个分量从物理意义明确的起点出发，以减少陷入局部最优的概率。完整流程如算法~\ref{alg:template_recover} 所示。

\begin{algorithm}[htbp]
  \caption{离线模板构建与在线密钥恢复}
  \label{alg:template_recover}
  \KwIn{已知密钥 $k_b^{\mathrm{known}}$，特征缓存行集合 $\mathcal{R}$，每索引采样数 $S$，
  在线观测集 $\{(p_b^{(i)}, \mathbf{x}^{(i)})\}_{i=1}^{M}$}
  \KwOut{密钥字节估计 $\hat{k}_b$}
  \BlankLine
  \tcp{信号采集与模板训练}
  确定特征缓存行延迟区间 $\{[d_j^{\mathrm{lo}}, d_j^{\mathrm{hi}})\}$ \;
  \For{$e \leftarrow 0$ \KwTo $255$}{
    令 $p_b \leftarrow e \oplus k_b^{\mathrm{known}}$，采集 $S$ 次特征向量集合 $X^{(e)}$\;
    $\mathrm{GMM}[e] \leftarrow \mathrm{EM}(X^{(e)},\;\text{分量数}=2,\;\text{对角协方差})$\;
  }
  \BlankLine
  \tcp{在线密钥恢复}
  \For{$k \leftarrow 0$ \KwTo $255$}{
    $S(k) \leftarrow \displaystyle\sum_{i=1}^{M} \log p\!\left(\mathbf{x}^{(i)} \mid \mathrm{GMM}[p_b^{(i)} \oplus k]\right)$\;
  }
  \KwRet{$\hat{k}_b \leftarrow \arg\max_k\, S(k)$}
\end{algorithm}

当离线模板构建完成后，在线恢复阶段不再需要已知密钥，攻击者以随机明文触发目标加密，采集特征向量 $\mathbf{x}^{(i)}$，对 256 个候选密钥值 $k$ 分别计算累积对数似然。

\begin{equation}
  S(k) = \sum_{i=1}^{M} \log p\!\left(\mathbf{x}^{(i)} \;\middle|\; \mathrm{GMM}\!\left[p_b^{(i)} \oplus k\right]\right)
  \label{eq:gmm_score}
\end{equation}

候选密钥字节 $k$ 会决定明文字节 $p_b$ 被映射到哪个查表索引，即 $e = p_b \oplus k$。在线恢复时，系统对每个候选值都按这一关系计算对应的索引，再将当前观测到的特征向量与该索引对应的高斯混合模型进行比较。若候选值等于真实密钥字节，则它给出的索引与来宾实际访问的表项一致，观测特征与对应模型更匹配，因此得到较大的对数似然。若候选值错误，索引会偏离真实访问位置，观测特征与模型的匹配程度下降，得到的对数似然也会较小。

系统对全部观测样本重复这一过程，并将每个候选值在所有样本上的对数似然累加，得到总分 $S(k)$。总分越大，说明该候选值越能解释观测到的信号。随着样本数增加，正确候选与错误候选之间的分数差通常会逐步拉大，因此正确密钥更容易排到前面。为避免两个高斯分量在数值计算中出现下溢，分量合并时采用 log-sum-exp 形式计算总对数似然。

上述流程对字节位置 0、4、8、12 分别独立执行，各位置均通过 Te0 的查表访问产生可观测信号，共享同一套特征构造与模型评分机制。四个字节位置的恢复结果合并后，可得到 AES-128 第一轮中与 Te0 查表对应的 4 个密钥字节，即 $k_0, k_4, k_8, k_{12}$，占完整 128 位密钥的四分之一。若 Te1、Te2、Te3 也存在跨页特征，则可继续采用相同方法分别定位各自的跨页结构并建立模板，则可覆盖第一轮全部 16 个密钥字节，从而完整恢复 AES-128 主密钥。

\section{AES T-Table 物理页定位}
\label{sec:exp_phase1}

第一阶段的目标是在不读取来宾内存内容的前提下，仅凭宿主机对 NPT 页表的观测和跨域一致性探测，确定 \texttt{Te0} 所在物理页的基地址 $\mathit{te0\_page\_gpa}$ 及其页内偏移 $\mathit{te0\_inpage\_off}$，为后续模板构建提供所需的地址信息。定位过程分为两步：第一步利用 NPT Accessed 位差分扫描在全 GPA 空间中筛出与 AES 执行相关的候选页簇；第二步对候选页簇施加一致性信号验证，从多个高分候选中区分出真实的 Te0 页。

\subsection{静态页内偏移提取}

$\mathit{te0\_inpage\_off}$ 可在攻击发起前离线提取，不依赖来宾运行时状态。攻击者持有与来宾运行时相同版本的目标二进制，使用符号表工具读取 \texttt{Te0} 的虚拟地址 $\mathit{te0\_vma}$，再结合其所在 ELF LOAD 段的虚拟起始地址和文件偏移计算出文件内偏移 $\mathit{te0\_file\_off}$，根据公式\eqref{eq:inpage_off}对页大小取模。

在本实验所用的 OpenSSL 3.4.0 \texttt{no-asm} 版本中，$\mathit{te0\_vma} = \texttt{0x409f60}$，$\mathit{te0\_file\_off} = \texttt{0x9f60}$，由此得到 $\mathit{te0\_inpage\_off} = \texttt{0xf60}$。Te0 数组每个表项占 4 字节，从该页内偏移算起，前 $(4096 - 3936) / 4 = 40$ 个表项落在当前物理页内，即 Te0[0..39]，其余 Te0[40..255] 跨入下一物理页，由此得到跨页分界索引 $e^* = 39$。来宾以 \texttt{nokaslr norandmaps} 参数启动，用户态虚拟地址布局固定，因此 $\mathit{te0\_inpage\_off}$ 在每次虚拟机启动后保持不变，无需重复提取。$\mathit{te0\_page\_gpa}$ 则会随每次启动的物理内存分配变化，必须在线获取。

\subsection{NPT Accessed 位差分扫描}

AMD 的 NPT 页表项中包含一个由 MMU 硬件维护的 Accessed 位，来宾 vCPU 访问对应物理页时该位自动置 1。宿主机通过 KVM 内核扩展接口可以批量清零指定 GPA 范围内的 Accessed 位，并在稍后重新扫描哪些页面在该窗口内被访问过。差分扫描每轮分基线和触发两个阶段。基线阶段在清零 Accessed 位后等待约 5 毫秒，期间内核调度、内存管理等系统活动正常进行，其产生的页面访问记录为各页的 $\mathrm{base}[p]$ 计数。触发阶段重新清零后立即向来宾 AES 服务连续发送 50 次随机明文请求，将观测窗口限定在加密执行期间，扫描结果记录为 $\mathrm{trig}[p]$。两阶段的差值 $\mathrm{score}(p) = \mathrm{trig}(p) - \mathrm{base}(p)$（见式~\eqref{eq:npt_score}）衡量页面在 AES 加密执行期间相对于背景的额外访问频率。实验共进行 100 轮，扫描范围为 \texttt{0x0} 至 \texttt{0x200000000}。

扫描范围需同时覆盖低地址段和 4GB 以上的私有页别名区。在 SEV-SNP 环境下，QEMU 通过 \texttt{guest\_memfd} 机制将来宾私有内存单独管理，这部分内存在宿主 GPA 视图中会在 4GB 以上形成一段别名映射窗口。实验中观测到真实的 $\mathit{te0\_gpa}$ 落在 \texttt{0x101000000} 以上的高地址区域，因此扫描终止地址取低地址段上界与私有别名段上界中的较大值，以保证目标页落在扫描范围内。

差分扫描结束后，大量页面的 $\mathrm{score}$ 均达到满分（$\mathrm{trig} = 100$、$\mathrm{base} = 0$）。AES 执行期间除访问 T-Table 外，还会访问代码段、程序栈、动态链接库数据以及网络收发缓冲区，这些页面在触发阶段同样被频繁访问，因此单纯依据差分得分无法从数百个高分页面中定位目标。对此，算法将高分候选页按物理地址邻近性（相邻间距不超过 2 页）归并为页簇，再以总差分得分之和、页数和地址跨度三元组对页簇排序，总得分越高、页数越少、跨度越小的页簇排名越靠前（见式~\eqref{eq:cluster_score}）。Te0--Te3 四张表的物理页在进程地址空间中连续排列，对应的候选页簇具有得分高、页数适中、跨度紧凑的特征，在聚类排序中处于靠前位置。取得分最高的页簇作为怀疑区间，再向两端各扩展 2 页，最终将候选范围从约 52 万页收窄至 5 页左右。

\subsection{一致性信号多特征验证}

通过差分扫描给出怀疑区间后，仍需进一步确认哪个候选页是真实的 Te0 页。实验表明，差分扫描选出的得分最高页簇并不总是包含 Te0 的真实页簇，页簇得分排名与实际目标并不严格对应。因此需要第二步对候选页簇中的每个代表页进行一致性信号验证，以区分真实的 T-Table 页面与其他高分竞争页面。

验证过程分两个维度。其一是输入相关性检验，攻击者控制明文，使 Te0 的查表索引 $e_0 = p_0 \oplus k_0$ 在候选页覆盖的区间内系统性变化，观察宿主侧对该候选页的缓存行延迟是否随 $e_0$ 的变化而呈现出规律性的分布差异。真实 Te0 页在来宾执行查表时会产生跨域一致性逐出，宿主侧延迟随 $e_0$ 的变化而变化；而无关页面则不会对输入产生可重复的响应，延迟分布在不同 $e_0$ 之间近似一致，相关性分数 $\rho$ 偏低。其二是一致性信号信噪比：对候选页上的每条缓存行，在来宾未访问该页（H0）和来宾查表触及该页（H1）状态下分别采集延迟分布，计算信噪比 SNR，取各行最大值作为该候选页的 $\mathrm{SNR}_{\max}$。SNR 衡量的是一致性逐出信号在该页上是否稳定可观测；对于真实 Te0 页，来宾查表触发跨域逐出的信号清晰，SNR 远高于不含 T-Table 的页面。

\begin{equation}
  C(i) = \bigl(\min(\mathrm{SNR}_{\max}^{(i)},\, 100) \times 0.3 + \rho^{(i)} \times 0.4\bigr) \times \alpha_i + B_i,
  \label{eq:combined_score}
\end{equation}

综合评分根据公式~\eqref{eq:combined_score}将两个维度的特征加权合并。其中 $\alpha_i$ 为页簇大小惩罚因子，簇规模越大则权重越低，以排除代码段等大型映射区域；$B_i$ 为基于 NPT 一致性特征的加分项，NPT 一致性加分考察两类辅助指标。第一类是簇内各页的基线命中次数是否一致，$\mathrm{base}[p]$ 值在簇内相同说明该簇的背景访问模式在多轮扫描中稳定。第二类是基线命中总次数是否处于合理范围，Te0 页所在区域在来宾 AES 服务持续运行期间，因预热调用等原因，基线阶段也会产生少量访问，$\mathrm{base}$ 值略大于零，这与代码段等纯触发期访问的页面有所区别。这两类特征在实验中对区分 Te0 页与代码段、栈等高分竞争页面起到了辅助作用。表~\ref{tab:phase1_cluster_ranking} 给出一次典型实验中差分扫描产生的候选页簇经一致性验证后的排名结果，包含 5 个得分较高的候选簇。

\begin{table}[htbp]
  \centering
  \caption{候选页簇综合评分排名}
  \label{tab:phase1_cluster_ranking}
  \begin{tabular}{cccccc}
    \toprule
    排名 & 页簇代表地址 & $\mathrm{SNR}_{\max}$ & 相关性 $\rho$ & 综合评分 & 是否目标 \\
    \midrule
    1 & \texttt{0x1ebe000}    & 42.93 & 7.31  & 73.68 & 是 \\
    2 & \texttt{0x10893c000}  &  6.81 & 6.99  & 71.88 & 否 \\
    3 & \texttt{0x101c5b000}  &  4.85 & 6.17  & 70.38 & 否 \\
    4 & \texttt{0x2136000}    & 32.92 & 5.00  & 66.76 & 否 \\
    5 & \texttt{0x3355000}    & 21.15 & 12.58 & 62.72 & 否 \\
    \bottomrule
  \end{tabular}
\end{table}

图~\ref{fig:cluster_ranking} 和表~\ref{tab:phase1_cluster_ranking} 共同说明了引入一致性信号的必要性。从综合评分来看，目标页簇排名第 1（73.68），次优候选（71.88）与之仅相差 1.8 分，单凭综合评分的绝对差值不足以保证区分可靠。但从 $\mathrm{SNR}_{\max}$ 来看，目标页簇（42.93）是排名第 2 候选簇（6.81）的 6.3 倍，信号区分度远高于仅有略微差异的综合评分所能反映的程度。第 4 位候选簇虽然 SNR 也较高（32.92），但其 NPT 一致性加分为零（不满足基线一致性条件），综合评分因此下降至第 4 位。这说明单一特征均不足以可靠区分目标页，几类特征配合使用才能有效识别真实 Te0 页。

在本文进行的多次实验中，目标 Te0 页在综合评分排名中均位于前三位，其中实验多数情况位于第 1 位。结合对前三名候选页簇分别执行短暂的模板匹配验证，可以在后续阶段进一步确认真实目标页。定位阶段的主要贡献是将候选范围从约 52 万页压缩至不超过 5 页的候选集，并给出一个可靠的优先排序，使第二阶段的模板构建代价控制在可接受范围内。

\begin{figure}[htbp]
  \vspace{-0.3em}
  \centering
  \includegraphics[width=0.92\textwidth]{pic/cluster_ranking.png}
  \caption{候选页簇综合评分与一致性信号信噪比对比}
  \label{fig:cluster_ranking}
\end{figure}

\section{AES 部分密钥恢复}
\label{sec:exp_phase2}

本阶段以第一阶段确定的 $\mathit{te0_page_gpa}$ 和 $\mathit{te0_inpage_off}$ 为输入，分离线模板训练和在线密钥恢复两步执行。实际实验表明，两类可观测信号在攻击流程中的作用并不相同。跨域一致性逐出信号能够稳定反映目标页及其邻近缓存行的结构信息，因此更适合用于目标页确认、边界附近缓存行筛选以及启动内的特征定位；但在单字节恢复阶段，该信号对不同查表索引的区分能力不够稳定。相比之下，基于 NOCACHE 探测得到的物理内存竞争信号在重复实验中表现出更好的判别性。特别是在将一次突发探测中的原始延迟序列转换为延迟区间计数特征后，不同查表索引在若干模板特征缓存行上的计数分布呈现出较稳定的差异。因此，当前密钥恢复阶段的最终模型仅使用竞争信号构造特征，而将一致性信号保留为页面定位和特征缓存行选择的辅助信息。

\subsection{实验流程}

离线训练在已知密钥的受控环境中进行。攻击者以已知密钥 $k_0^{\mathrm{known}}$ 控制明文字节 $p_0$，令查表索引 $e_0 = p_0 \oplus k_0^{\mathrm{known}}$ 遍历 0 到 255，对每个索引值重复触发加密并采集竞争信号。每次触发后，宿主机立即对一组预先选定的模板特征缓存行执行 NOCACHE 突发探测，每条缓存行连续探测 16 次，记录各次延迟读数作为一条原始观测的原始序列。将原始序列压缩为延迟区间计数特征向量（详见式~\eqref{eq:feature_component}）后，每条观测转化为一个 11 维特征向量。每个索引的特征向量集合随后用于拟合两分量对角高斯混合模型（详见式~\eqref{eq:gmm_model}），具体流程见算法~\ref{alg:template_recover}。

本实验共进行 16 轮离线采集，每轮对 256 个索引各收集 16 条样本，16 轮数据合并后统一训练，每个索引的训练样本量约为 256 条。训练数据包含同一 VM 会话内的多次重复，能够覆盖竞争信号在不同采集时刻下的正常波动。

在线恢复阶段以未知密钥触发加密，采集与训练阶段格式相同的特征向量序列，对 256 个候选密钥字节 $k$ 分别累加对数似然（见式~\eqref{eq:gmm_score}），取总分最高者作为密钥字节估计 $\hat{k}_0$。

\subsection{模板特征缓存行的选取}

本实验为字节位置 0 选取了两类模板特征缓存行，共 9 条基础缓存行加 2 个派生维度，合计 11 维特征向量。

第一类特征是位于 Te0 第二物理页低地址区域的两条缓存行。这两条缓存行在物理上紧邻 Te0 的跨页边界，所处 DRAM bank 与 Te0[0..39] 所在第一物理页的部分区域存在交叉响应关系。对这两条缓存行的 NOCACHE 探测延迟在不同索引下呈现出明显的分布差异，Fisher 准则分数在所有特征中位列最高，成为当前模型的主要判别特征。

第二类特征是对 Te0 第一物理页执行全页缓存行延迟扫描后，按一致性信号信噪比排名靠前的 7 条缓存行。这些缓存行位于 Te0 的主存储页上，其 NOCACHE 探测延迟在不同索引下的变化规律来自 DRAM Interleaving 带来的 bank 争用模式，为第一类特征无法覆盖的索引区间提供补充判别能力。

此外，将第一类两条缓存行的计数值取最大值和求和，分别作为两个派生维度加入特征向量，以增强对两条缓存行共同响应的描述。

\subsection{延迟区间计数特征与 d-range 选取}

对每条模板特征缓存行，NOCACHE 突发探测的原始输出是一组长度为 16 的延迟读数序列。早期实验直接对序列求均值作为特征值，但均值对调度抖动和个别极高延迟读数敏感，相同索引下不同轮次的均值波动较大，导致同一索引的特征样本分散，不同索引之间的均值差异被噪声淹没。

为解决这一问题，将每条缓存行的特征改为延迟区间计数：预先为该缓存行确定一个延迟区间 $[d_j^{\mathrm{lo}},\, d_j^{\mathrm{hi}})$，统计突发序列中落入该区间的次数作为特征值（见式~\eqref{eq:feature_component}）。这一表示把"平均延迟偏移了多少"转化为"竞争事件在窗口内出现了多少次"，极端噪声值只影响单次读数是否落入区间，而不会拉高整个均值，特征值的稳定性因此提高。

每条缓存行的延迟区间 $[d_j^{\mathrm{lo}},\, d_j^{\mathrm{hi}})$ 称为该特征的 d-range，由训练数据自动搜索确定，而不是人工设定固定阈值。搜索时遍历候选下界 $d^{\mathrm{lo}}$，对每个候选区间计算 Fisher 准则分数：

\begin{equation}
  F_j = \frac{s_{\mathrm{between}}^2}{s_{\mathrm{within}}^2}
  \label{eq:fisher_criterion}
\end{equation}

其中 $s_{\mathrm{between}}^2$ 为各索引 $e$ 对应计数均值之间的方差，衡量该区间下不同索引的平均计数值差异；$s_{\mathrm{within}}^2$ 为同一索引内不同样本计数值之间的方差均值，衡量同一索引下特征值的波动。Fisher 分数越高，说明该区间在不同索引之间的区分能力越强，同时对同一索引的观测越稳定。取分数最高的候选区间作为最终 d-range，同时避免了区间过宽时计数长期饱和（所有读数均落入区间导致不同索引的计数无差异）以及区间过窄时有效事件被遗漏的问题。

本实验 byte 0 模型的 d-range 值如表~\ref{tab:drange_byte0} 所示。

\begin{table}[htbp]
  \centering
  \caption{byte 0 模型各特征的延迟区间（d-range）与 Fisher 分数}
  \label{tab:drange_byte0}
  \begin{tabular}{lccc}
    \toprule
    特征描述 & $d^{\mathrm{lo}}$（周期） & $d^{\mathrm{hi}}$（周期） & Fisher 分数 \\
    \midrule
    第二物理页缓存行 1  & 144 & 400 & 1.537 \\
    第二物理页缓存行 2  & 120 & 376 & 1.536 \\
    第一物理页高响应行 1 & 120 & 248 & 0.684 \\
    第一物理页高响应行 2 & 120 & 376 & 1.180 \\
    第一物理页高响应行 3 & 120 & 248 & 0.653 \\
    第一物理页高响应行 4 & 408 & 664 & 0.250 \\
    第一物理页高响应行 5 & 424 & 680 & 0.395 \\
    第一物理页高响应行 6 & 120 & 248 & 1.544 \\
    第一物理页高响应行 7 & 120 & 376 & 0.790 \\
    \bottomrule
  \end{tabular}
\end{table}

两条第二物理页缓存行的 Fisher 分数均在 1.53 以上，是所有特征中最高的，说明这两条缓存行的区间计数在 256 个索引之间的区分能力最强。第一物理页的 7 条高响应缓存行中，高响应行 4 和高响应行 5 的 d-range 明显偏高（下界 408--424 周期），其余 5 条的下界在 120 周期附近，说明这两类缓存行所反映的竞争延迟范围不同，对应 DRAM 不同层级的访问争用。

\subsection{特征模式的结构分析}

图~\ref{fig:gmm_template_heatmap} 展示了训练所得模型中，256 个查表索引对应的高计数分量均值在 9 个基础特征维度上的分布，每列按最大值归一化。从热力图中可以观察到，after$_0$ 和 after$_1$ 呈现出约 32 个索引为一组的交替条带，这一周期与 DRAM Interleaving 的物理地址映射一致。table\_top$_4$ 和 table\_top$_5$ 的条带模式与 after 特征大体互补：after 组计数较高时 top$_4$ 和 top$_5$ 较低，反之亦然。这表明两类特征覆盖了来宾访问 DRAM 不同物理位置时在宿主侧产生的互补响应，共同为 256 个索引提供可区分的特征签名。

\begin{figure}[htbp]
  \centering
  \includegraphics[width=\textwidth]{pic/gmm_template_heatmap.png}
  \caption{byte 0 GMM 模板的特征均值热力图。左图为 9 维特征在 256 个查表索引上的归一化高计数分量均值；右图为边界特征 after$_0$ 与高响应特征 table\_top$_4$ 的互补模式。}
  \label{fig:gmm_template_heatmap}
\end{figure}

\subsection{密钥恢复效果}

实验以字节位置 0 的离线模板对 24 轮独立在线采样执行打分，统计正确密钥字节（$k_0 = \texttt{0x2b}$）在 256 个候选中的排名分布。结果见表~\ref{tab:partial_recovery_results}。

\begin{table}[htbp]
  \centering
  \caption{byte 0 密钥字节恢复情况}
  \label{tab:partial_recovery_results}
  \begin{tabular}{lccccc}
    \toprule
    轮数 & Top-1 & Top-3 & Top-5 & Top-8 & 平均排名 \\
    \midrule
    24 & 17/24 (70.83\%) & 19/24 (79.17\%) & 21/24 (87.50\%) & 24/24 (100\%) & 2.25 \\
    \bottomrule
  \end{tabular}
\end{table}

24 轮密钥恢复实验中，正确密钥字节 Top-1 成功率为 70.83\%，Top-8 达到 100\%，平均排名 2.25。结果表明，当前离线模板在独立在线采样条件下已能够将真实密钥稳定约束到少量候选之中，增加采样轮次或引入自适应校准有望进一步提升 Top-1 成功率。
